\subsection{Conclusion}
This thesis has successfully resolved the critical ``Master Gap'' identified in contemporary smart grid literature by constructing a unified, Tri-layer Cooperative Resilience Architecture. By mathematically synthesizing Second-Order Cone Programming (SOCP) for exact physical bounding, Analytical Target Cascading (ATC) for secure spatial decentralization, and Model Predictive Control (MPC) for proactive temporal dynamism, this research has developed a fault-adaptive, P2P energy trading framework capable of ensuring extreme Microgrid resilience. The efficacy of this unified architecture is summarized by four key findings:
\begin{itemize}
    \item \textbf{Algorithmic Architecture:} The proposed Bottom-up Distributed Dynamic 3-Mode State Machine achieved high computational tractability, reaching rapid ATC convergence (max 91 seconds) during severe structural grid shocks to ensure real-time MPC feasibility.
    \item \textbf{Physical Resilience and AC-OPF Necessity:} The absolute necessity of AC-OPF was validated by demonstrating that proactive active power curtailment effectively mitigates catastrophic voltage collapse caused by transit-induced reactive losses, strictly maintaining the 0.90~p.u. boundary.
    \item \textbf{Market Pricing and the Tie-line Congestion Paradox:} The study mathematically proved the Tie-line Congestion Paradox, wherein physical thermal limits inherently activate KKT bounds to forcefully cap ATC consensus pricing ($\lambda$), completely preventing market collapse during extreme localized VOLL escalations.
    \item \textbf{System Security and Graceful Degradation:} The framework synergizes a Graceful Degradation mechanism with P2P trading to guarantee a 100\% Critical Load survival rate while achieving a Negative Premium---a dual economic-resilience benefit that liberates stranded energy during severe disruptions.
\end{itemize}

Ultimately, this thesis fundamentally establishes a critical paradigm shift in decentralized energy management. It mathematically proves that P2P trading platforms must transcend their conventional role as mere financial markets designed for economic optimization. Instead, they can be weaponized as a structural, physical safety net, seamlessly harmonizing economic profitability under normal conditions with catastrophic fault survivability during extreme grid emergencies.

\subsection{Future Work}
To fundamentally address the physical and algorithmic limitations observed in this study, future research should focus on the following primary directions:
\begin{itemize}
    \item \textbf{Hardware-in-the-Loop (HIL) Real-Time Validation:} While the proposed Tri-layer Architecture demonstrated exceptional resilience in simulation, deploying the algorithms into physical microcontrollers via HIL testbeds is a mandatory next step to evaluate real-world communication latencies, sensor noise, and hardware-execution delays.
    \item \textbf{Asynchronous ATC for Massive Scalability:} To scale the framework from a small cluster of MGs to city-wide smart grids (100+ MGs), future iterations must transition from synchronous ATC to asynchronous or event-triggered consensus algorithms to mitigate communication bottlenecks and ensure the MPC rolling horizon remains uncompromised.
    \item \textbf{Overcoming Conservative BESS SOC Retention:} To mitigate the conservative MPC BESS discharging behavior (driven by the strict end-of-horizon SOC penalty constraints), introducing a Reinforcement Learning (RL)-based adaptive SOC Penalty will encourage more proactive BESS utilization. Additionally, future iterations will employ Stochastic MPC to better handle the inherent uncertainty of renewable energy forecasting.
    \item \textbf{Hardware and Reactive Power Support:} Currently, when handling massive transit currents that threaten voltage stability, intermediary MGs are forced to suboptimally shed active power. To resolve this limitation, future iterations will integrate dynamic reactive compensation devices (e.g., STATCOM, SVC) or Smart Inverters equipped with Volt-VAR control to independently supply the necessary reactive power ($Q$) support.
    \item \textbf{Dynamic Topology Reconfiguration:} To autonomously isolate physical fault zones and establish backup routing paths, future work will integrate a Dynamic Switch Control mechanism for real-time network reconfiguration, thereby optimizing power flow and further enhancing the structural resilience of the MMG network.
\end{itemize}
